\section{Related Work}
We discuss recent studies on self-supervised speech representation learning by grouping them by training objective. The earliest line of work learns representations by postulating a generative model for speech with latent variables, which are assumed to capture the relevant phonetic information. Training of these models amounts to likelihood maximization. Different latent structures have been applied to encode the prior assumption, such as continuous~\cite{hsu2017learning}, discrete~\cite{chorowski2019unsupervised,van2017neural}, or sequential~\cite{hsu2017unsupervised,ebbers2017hidden,glarner2018full,khurana2019factorial,khurana2020convolutional}. 


Prediction-based self-supervised learning has gathered increasing interests recently, where a model is tasked to predict the content of the unseen regions \cite{chung2019unsupervised, chung2020generative, chung2020improved, ling2020deep, wang2020unsupervised, liu2020mockingjay, chi2020audio, ling2020decoar} or to contrast the target unseen frame with randomly sampled ones \cite{oord2018representation, kharitonov2020data, schneider2019wav2vec, baevski2020wav2vec}. Some models combine both the predictive and the contrastive losses \cite{baevski2019vq, baevski2019effectiveness}. These objectives can usually be interpreted as mutual information maximization~\cite{tsai2020ssl_multi}. Other objectives do not belong to these categories, for example, \cite{pascual2019learning}.

This work is most related to DiscreteBERT~\cite{baevski2019effectiveness}: both HuBERT and DiscreteBERT predict discrete targets of masked regions. However, there are several crucial differences. First, instead of taking quantized units as input, HuBERT takes raw waveforms as input to pass as much information as possible to the transformer layers, which was shown to be important in \cite{baevski2020wav2vec}. Furthermore, in the experiment section, we show that our model, with simple k-means targets, can achieve better performance than DiscreteBERT that uses vq-wav2vec \cite{baevski2019vq} learned units. Second, we also present many techniques to improve teacher quality instead of using a single fixed teacher as done in DiscreteBERT.


HuBERT is also related to wav2vec 2.0~\cite{baevski2020wav2vec}. However, the latter employs a contrastive loss that requires careful design of where to sample negative frames from, an auxiliary diversity loss to encourage the discrete unit usage, and demands a proper Gumbel-softmax temperature annealing schedule. In addition, it only explores quantizing the waveform encoder output, which may not be the best feature for quantization due to the limited capacity of the convolutional encoder, as suggested by our ablation studies in Figure~\ref{fig:qual_layer}. Concretely, our proposed method adopts a more direct predictive loss by separating the acoustic unit discovery step from the masked prediction representation learning phase and achieves the state-of-the-art results that match or outperform wav2vec 2.0 on different fine-tuning scales.

Finally, the idea of iterative refinement target labels is similar to iterative pseudo labeling for semi-supervised ASR~\cite{xu2020iterative, likhomanenko2020slimipl}, which leverages an improving student model to generate better pseudo-labels for the next iteration of training. The HuBERT approach can be seen as extending this method to the self-supervised setup with a masked prediction loss.